\section*{ОБОЗНАЧЕНИЯ И СОКРАЩЕНИЯ}

БШ -- блочное шифрование.

ИТ -- информационные технологии. 

ПО -- программное обеспечение.

ТЗ -- техническое задание.

ТП -- технический проект.

AES (Advanced Encryption Standard) -- симметричный блочный алгоритм шифрования.

API (Application Programming Interface) -- программный интерфейс приложения, набор готовых функций для использования в сторонних модулях.

ECB (Electronic Codebook) -- режим электронной кодовой книги, режим шифрования, при котором каждый блок данных шифруется независимо.

PKCS\#7 (Public-Key Cryptography Standards \#7) -- стандарт дополнения данных до кратности размеру блока.

S-box (Substitution-box) -- таблица замен, используемая в алгоритмах симметричного шифрования.

SBCL (Steel Bank Common Lisp) -- компилятор и среда выполнения языка Common Lisp.